\section{Hilo conductor de la clase: la tijera entre el aumento de capacidades y la disminución de la apertura}

\subsection{Juicio inicial: Why Access Matters}
Percy Liang plantea, en los primeros minutos, la tesis central de toda la clase: \textbf{las capacidades de los Foundation Model aumentan rápidamente, pero el nivel de acceso disponible para los investigadores disminuye en paralelo}. Esta tijera de «curva de capacidad hacia arriba, curva de apertura hacia abajo» no es solo un cambio de estrategia comercial, sino que reconfigurará directamente la agenda de investigación de los próximos cinco años.

\begin{importantbox}{Access Shapes Research}
El expositor lo resume en una frase: \textit{``Access shapes research.''}
Cuando solo se tiene acceso por API, lo que más fácilmente se puede hacer es Prompt Engineering; cuando se cuenta con los pesos y el código de entrenamiento, se abre la posibilidad de hacer interpretabilidad mecanicista, modificación del entrenamiento, auditoría de seguridad e innovación arquitectónica.
\end{importantbox}

\begin{figure}[H]
\centering
\includegraphics[width=0.93\textwidth]{frame-01-title.jpg}
\caption{Página de apertura de la clase: del tema de los Agent hacia una pregunta abierta de ciencia}
\end{figure}
\noindent\footnotesize Evidencia visual: los subtítulos se ubican aproximadamente entre 00:00:09--00:00:51, donde el expositor plantea explícitamente la ``importance of open source'' y ``access shapes research''.\normalsize

\subsection{Perspectiva histórica: los tres “dividendos de acceso” del NLP}
El expositor repasa varios saltos característicos en la historia del NLP y el aprendizaje automático:
\begin{enumerate}
    \item de los años noventa a los dosmiles: la mayor accesibilidad de texto en internet impulsó el NLP estadístico;
    \item la década de 2010: la maduración de las plataformas de etiquetado de datos impulsó el aprendizaje supervisado y la evaluación mediante benchmarks;
    \item la era del aprendizaje profundo: la mayor disponibilidad de GPU y clústeres de cómputo hizo posible escalar el transformer.
\end{enumerate}

\begin{knowledgebox}{Metodología implícita a nivel de curso}
Este repaso histórico no es un «paréntesis de divulgación», sino que prepara una metodología para lo que sigue:
\textbf{al discutir las capacidades de un modelo, es necesario discutir a la vez el acceso a los datos, el acceso al cómputo y el acceso al sistema.}
\end{knowledgebox}

\begin{warningbox}{Error común: reducir la pregunta de apertura a «si los pesos son de código abierto»}
La pregunta de apertura no es un interruptor binario. Que los pesos puedan descargarse es solo una dimensión. Si los datos de entrenamiento, el proceso de entrenamiento, los scripts de evaluación, los registros del sistema y la estrategia de despliegue no son visibles, la reproducibilidad de la investigación sigue siendo insuficiente y las conclusiones científicas siguen siendo frágiles.
\end{warningbox}

\subsection{Resumen de la sección}
Esta sección plantea la pregunta central de toda la clase: no se trata de debatir si «el código abierto es bueno o no», sino de discutir qué investigación permite y qué investigación bloquea cada nivel de acceso, y cómo esta diferencia estructural afecta el progreso a largo plazo de los Agent y los Foundation Model.

